\section{Front-camera videos of one seed under four weather conditions}\label{sec:videos}

\subsection{What was recorded and why}
All 146 earlier videos (about 1.2\,GB: top-down and bird's-eye-view panels, laggy camera panels, and a front-camera set at a different seed with 13 models) were deleted. They are replaced by one small, consistent set, as requested: \textbf{front-camera frames only, one seed, three models, four weather conditions} (12 videos, \pth{report/videos/}).
The camera is the front-left camera of the teammates' rig (1200$\times$900 pixels) with the teammates' weather augmentation applied to each frame, written at 10 frames per second, \emph{one frame per simulation step} (so playback is in real time and smooth), with no text, boxes or overlays.

\paragraph{The three models:} the baseline behaviour cloning trained without traffic, the baseline behaviour cloning trained with sparse traffic, and the headline model.
\paragraph{The four weather conditions:} clear, fog 1.0, rain 0.6, rain 1.0 (the cells of the 200-episode evaluation).

\subsection{How the seed was chosen}
The seed was searched in the 200-episode result files (\pth{bev_rl/adaptive/results/n200_all_weather/} and \pth{final_200seed/}) for episodes in which \emph{both baselines fail and the headline model succeeds}, preferring seeds where this holds in all four weather cells.
Seven seeds satisfy the clear-weather condition with two collisions (5006, 5060, 5102, 5119, 5161, 5169, 5181). Of these, only \textbf{seed 5169} has the sparse-traffic baseline crashing in all four weather conditions \emph{and} the headline model succeeding in all four (the no-traffic baseline was evaluated in clear weather only at this sample size, where it crashes as well).
The videos were then recorded and their outcomes compared with the evaluation: in all conditions recorded the outcome is identical to the evaluation's and the episode length differs by at most a few steps (Section~\ref{sec:repro} explains why exact equality is not guaranteed).

\subsection{Result}
\begin{table}[H]\centering\small
\caption{The twelve recorded videos (seed 5169, sparse traffic, front-left camera). File names in \code{report/videos/} follow the pattern \code{front_camera__seed5169__<model>__<weather>__<outcome>.mp4}.}\label{tab:videos}
\resizebox{\textwidth}{!}{\input{tables/video_outcomes}}
\end{table}

\begin{figure}[p]\centering
\includegraphics[width=\textwidth]{figures/fig_seed5169_front_camera_frames.pdf}
\caption{Frames of the recorded videos (seed 5169). Each row is one model in one weather condition; columns are moments of the episode (the baselines crash at about step 247--264, i.e.\ 25 seconds; the headline model drives on until it arrives at about step 700). Rain 1.0 visibly degrades the image, yet the behaviour of all three models is the same as in clear weather.}
\label{fig:frames}
\end{figure}

\subsection{How to read the videos correctly}
\begin{itemize}
\item \textbf{The policies do not read these pixels.} They drive from the simulator-derived grid (Section~\ref{sec:surrogate}), with weather acting through detector dropout / calibrated degradation. The video shows what the real camera would be looking at under the same weather level next to the behaviour of the policy.
  That the baselines crash and the headline model does not, identically in every weather condition, \emph{is the weather result of Section~\ref{sec:weather} in video form}: the weather in the picture does not change what the models do.
\item The weather augmentation draws fresh random noise at every frame (as in the teammates' pipeline), so the rain texture flickers; this is expected.
\item The videos are not part of any statistic. They illustrate one scenario that was selected \emph{because} the baselines fail on it; they must not be used to estimate how often that happens (the 200-episode result does).
\end{itemize}

\subsection{How to regenerate}
\begin{lstlisting}
venv\Scripts\python.exe -m bev_rl.adaptive.record_camera_plain --out report\videos_raw --seeds 5169 ^
    --arms repo_bc_clear repo_bc_sparse ov_split --weather clear:0.0 fog:1.0 rain:0.6 rain:1.0
venv\Scripts\python.exe report\make_video_assets.py --finalise
\end{lstlisting}
Recording takes about 50 s per baseline video and 2.5 minutes per headline-model video (longer episode).
